\documentclass[11pt]{article}
\usepackage[margin=1in]{geometry}
\usepackage{amsmath,amssymb}
\usepackage[hidelinks]{hyperref}

\newcommand{\dd}{\mathrm{d}}
\newcommand{\ee}{\mathrm{e}}
\newcommand{\kB}{k_{\mathrm B}}
\newcommand{\Lloc}{\mathcal{L}}

\title{Derivations for ``Thermoelectric transport in hexagonal CrN nanoribbons''\\
\large Complete, step-by-step working of every formula used in the code and manuscript}
\author{}
\date{}

\begin{document}
\maketitle
\noindent
This companion note derives, without skipping steps, every mathematical result used in the
tight-binding + Landauer thermoelectric calculation: the lattice/reciprocal geometry, the
Slater--Koster Bloch Hamiltonian and its selection rules, the Onsager transport coefficients from
the Landauer currents (including all unit factors), the electronic figure-of-merit identity, the
Sommerfeld expansion with the Mott and Wiedemann--Franz consequences, the two-spin-channel
combination (with the bipolar correction), the ballistic phonon thermal-conductance quantum, and
the self-consistent mean-field edge magnetism with the Lichtenstein (LKAG) exchange formula.
Numerical cross-checks that were run against these results are noted in boxes.

\tableofcontents

%======================================================================
\section{Geometry and reciprocal lattice}

The honeycomb monolayer has two triangular sublattices, Cr ($A$) and N ($B$). Take primitive
vectors and nearest-neighbour (NN, $A\!\to\!B$) vectors
\begin{equation}
\mathbf a_1 = d\Big(\tfrac32,\tfrac{\sqrt3}{2}\Big),\quad
\mathbf a_2 = d\Big(\tfrac32,-\tfrac{\sqrt3}{2}\Big),\qquad
\boldsymbol\delta_1 = d(1,0),\ \
\boldsymbol\delta_2 = d\big(-\tfrac12,\tfrac{\sqrt3}{2}\big),\ \
\boldsymbol\delta_3 = d\big(-\tfrac12,-\tfrac{\sqrt3}{2}\big),
\end{equation}
where $d$ is the Cr--N bond length. Each is a unit-cell repeat of the others:
$\boldsymbol\delta_2=\boldsymbol\delta_1-\mathbf a_2$, $\boldsymbol\delta_3=\boldsymbol\delta_1-\mathbf a_1$,
and $\sum_i\boldsymbol\delta_i=0$.

\paragraph{Lattice constant.} The same-sublattice spacing is
$a\equiv|\mathbf a_1| = d\sqrt{(3/2)^2+(\sqrt3/2)^2} = d\sqrt{9/4+3/4}=d\sqrt3$, i.e.
\begin{equation}
\boxed{\,a=\sqrt3\,d\,}\qquad\Rightarrow\qquad
a=\sqrt3\times1.884~\text{\AA}=3.263~\text{\AA}\approx3.258~\text{\AA (DFT)}.
\end{equation}

\paragraph{Reciprocal vectors.} They satisfy $\mathbf a_i\!\cdot\!\mathbf b_j=2\pi\delta_{ij}$.
Writing $A=\left(\begin{smallmatrix}\mathbf a_1\\\mathbf a_2\end{smallmatrix}\right)$ (rows) and
$B=\left(\begin{smallmatrix}\mathbf b_1\\\mathbf b_2\end{smallmatrix}\right)$, the condition is
$A B^{\mathsf T}=2\pi\,\mathbb 1$, so $B=2\pi (A^{-1})^{\mathsf T}$. With $A=d\left(\begin{smallmatrix}
3/2 & \sqrt3/2\\ 3/2 & -\sqrt3/2\end{smallmatrix}\right)$, $\det A = d^2(3/2)(-\sqrt3)= -\tfrac{3\sqrt3}{2}d^2$,
\begin{equation}
\mathbf b_1=\frac{2\pi}{3d}\big(1,\sqrt3\big),\qquad
\mathbf b_2=\frac{2\pi}{3d}\big(1,-\sqrt3\big).
\end{equation}

\paragraph{Structure factor and Dirac points.} The $A\!\to\!B$ Bloch sum is
$f(\mathbf k)=\sum_{i}\ee^{\,i\mathbf k\cdot\boldsymbol\delta_i}$. It vanishes when the three unit
phasors sum to zero, i.e.\ when they are mutually $120^\circ$ apart. In fractional coordinates
$\mathbf k=k_1\mathbf b_1+k_2\mathbf b_2$ one finds $f=0$ at the zone corners
$(k_1,k_2)=(\tfrac13,-\tfrac13)$ and $(\tfrac23,\tfrac13)$ (the $K$, $K'$ points); at the zone
centre $|f|=3$.
\begin{quote}\small\itshape
Numerical check: evaluating $|f|$ at several high-symmetry fractional points confirmed
$|f(\tfrac23,\tfrac13)|=0$ while $|f(\tfrac13,\tfrac13)|=\sqrt3$; this fixed the correct $K$
label in the band-structure script.
\end{quote}

%======================================================================
\section{Tight-binding Bloch Hamiltonian}

\subsection{Orbital content and Bloch basis}
We keep the out-of-plane manifold that governs states near $E_F$: Cr $\{d_{z^2},d_{xz},d_{yz}\}$
on $A$ and N $\{p_z\}$ on $B$. Bloch states $|\mathbf k,\alpha,X\rangle
=\frac{1}{\sqrt{N_c}}\sum_{\mathbf R}\ee^{\,i\mathbf k\cdot(\mathbf R+\boldsymbol\tau_X)}
|\mathbf R,\alpha,X\rangle$ (sublattice $X\in\{A,B\}$, orbital $\alpha$) give a $4\times4$
matrix $H_{\alpha X,\beta Y}(\mathbf k)=\sum_{\mathbf R}\ee^{\,i\mathbf k\cdot(\mathbf R+\boldsymbol\tau_Y-\boldsymbol\tau_X)}
\langle \mathbf 0,\alpha,X|H|\mathbf R,\beta,Y\rangle$.

\subsection{Slater--Koster two-centre integrals for $p$--$d$ bonds}
For a bond with direction cosines $(l,m,n)$, the Slater--Koster rules give the $p$--$d$ matrix
elements in terms of two integrals $V_{pd\sigma},V_{pd\pi}$. The ones involving $p_z$ are
\begin{align}
E_{z,xy}&=\sqrt3\,l m n\,V_{pd\sigma}-2 l m n\,V_{pd\pi}, &
E_{z,yz}&=\sqrt3\,m n^2\,V_{pd\sigma}+m(1-2n^2)\,V_{pd\pi},\\
E_{z,zx}&=\sqrt3\,l n^2\,V_{pd\sigma}+l(1-2n^2)\,V_{pd\pi}, &
E_{z,x^2-y^2}&=\tfrac{\sqrt3}{2}\,n(l^2-m^2)V_{pd\sigma}-n(l^2-m^2)V_{pd\pi},\\
E_{z,3z^2-r^2}&=n\big[n^2-\tfrac12(l^2+m^2)\big]V_{pd\sigma}-\sqrt3\,n(l^2+m^2)V_{pd\pi}. &&
\end{align}
Here $E_{z,\alpha}\equiv\langle p_z|H|d_\alpha\rangle$ along the bond.

\subsection{Planar selection rule ($n=0$)}
The monolayer is planar (buckling $\pm0.07$~\AA\ is neglected), so every Cr--N bond lies in-plane:
$n=0$ and $l^2+m^2=1$. Substituting $n=0$ into the five expressions above,
\begin{equation}
E_{z,xy}=0,\quad E_{z,x^2-y^2}=0,\quad E_{z,3z^2-r^2}=0,\qquad
\boxed{\,E_{z,zx}=l\,V_{pd\pi},\quad E_{z,yz}=m\,V_{pd\pi}\,.}
\end{equation}
Thus \emph{N $p_z$ couples only to Cr $d_{xz}$ and $d_{yz}$} (through $V_{pd\pi}$), and
\emph{$d_{z^2}$ (and $d_{x^2-y^2}$) are non-bonding} with respect to $p_z$ for in-plane bonds. The
$d_{z^2}$ band therefore disperses only through Cr--Cr second-neighbour hopping.
\begin{quote}\small\itshape
Numerical check: for 20 random $\mathbf k$ and both spins, the code verified $H$ Hermitian and
$H_{d_{z^2},p_z}=0$ identically.
\end{quote}

\subsection{Bloch matrix elements}
Summing the three NN bonds (Cr at $\mathbf 0$, N at $\boldsymbol\delta_i$, cosines
$\hat{\boldsymbol\delta}_i=\boldsymbol\delta_i/d$),
\begin{equation}
f_{xz}(\mathbf k)=\sum_i \frac{\delta_{i,x}}{d}\,\ee^{\,i\mathbf k\cdot\boldsymbol\delta_i},\qquad
f_{yz}(\mathbf k)=\sum_i \frac{\delta_{i,y}}{d}\,\ee^{\,i\mathbf k\cdot\boldsymbol\delta_i},
\end{equation}
so $\langle d_{xz}|H|p_z\rangle=V_{pd\pi}f_{xz}$, $\langle d_{yz}|H|p_z\rangle=V_{pd\pi}f_{yz}$.
The Cr--Cr second-neighbour $d_{z^2}$--$d_{z^2}$ integral is direction independent at $n=0$
(from $E_{3z^2-r^2,3z^2-r^2}=\tfrac14 V_{dd\sigma}+\tfrac34 V_{dd\delta}\equiv t_{zz}$), giving a
triangular-lattice term $t_{zz}S_2(\mathbf k)$ with
$S_2(\mathbf k)=\sum_{\boldsymbol\tau}\ee^{\,i\mathbf k\cdot\boldsymbol\tau}$ over the six Cr--Cr
neighbours. Collecting terms with basis order $(d_{z^2},d_{xz},d_{yz},p_z)$ and adding the
mean-field exchange shift $\Delta_\sigma$ on the Cr $d$ orbitals,
\begin{equation}
H_\sigma(\mathbf k)=
\begin{pmatrix}
\varepsilon_{d_{z^2}}+\Delta_\sigma+t_{zz}S_2 & 0 & 0 & 0\\
0 & \varepsilon_\pi+\Delta_\sigma & 0 & V_{pd\pi}f_{xz}\\
0 & 0 & \varepsilon_\pi+\Delta_\sigma & V_{pd\pi}f_{yz}\\
0 & V_{pd\pi}f_{xz}^\ast & V_{pd\pi}f_{yz}^\ast & \varepsilon_{p_z}
\end{pmatrix}.
\end{equation}

\subsection{The effective conduction orbital $c$}
The four-orbital manifold above cannot host the majority \emph{conduction} band of the DFT
reference (electron pocket at $K$, minimum $-0.2$~eV; its projected weight lies in Cr
$d_{xy},d_{x^2-y^2},s$). We therefore add one effective orbital $c$ on the Cr triangular
sublattice. In a planar sheet an $s$- or in-plane-$d$-like Cr orbital has zero SK matrix element
with N $p_z$ at $n=0$ (the same table as above with $n=0$: $E_{s,z}=nV_{sp\sigma}=0$, and the
$E_{z,xy}, E_{z,x^2-y^2}$ rows vanish), so $c$ is symmetry-decoupled from the $\pi$ manifold and
enters $H$ as a diagonal block
\begin{equation}
H_{cc}(\mathbf k)=\varepsilon_c+\Delta^c_\sigma+\sum_{n=1}^{3}t_{cn}\,S_n(\mathbf k),\qquad
S_n(\mathbf k)=\sum_{\boldsymbol\tau\in\text{shell }n}\ee^{\,i\mathbf k\cdot\boldsymbol\tau},
\end{equation}
with shells at distances $a,\sqrt3 a,2a$ ($S_1(\Gamma){=}S_2(\Gamma){=}S_3(\Gamma){=}6$;
$S_1(K){=}-3$). The four constants $(\varepsilon_c,t_{c1},t_{c2},t_{c3})$ are a weighted least
squares fit to the digitized majority conduction band (weights favouring the transport window),
giving $(0.915,0.285,-0.028,0.030)$~eV: pocket at $K$ at $-0.20$~eV, saddle at $M$ at
$+0.58$~eV, $\Gamma$ at $+2.64$~eV (digitized: $-0.21$, $+0.63$, $+2.41$). In the ribbon
builders the three shells appear as the Cr--Cr \texttt{HoppingKind}s of range up to $2a$, which
forces a $2a$ lead period for the zigzag orientation.

\subsection{Mean-field (Stoner) exchange}
The on-site Hubbard term $U\,\hat n_{\alpha\uparrow}\hat n_{\alpha\downarrow}$ decoupled at
mean-field level gives $U\hat n_{\alpha\uparrow}\langle\hat n_{\alpha\downarrow}\rangle
+U\langle\hat n_{\alpha\uparrow}\rangle\hat n_{\alpha\downarrow}-U\langle\hat n_{\alpha\uparrow}
\rangle\langle\hat n_{\alpha\downarrow}\rangle$, so spin $\sigma$ Cr electrons see
$\varepsilon_{d}+U\langle n_{\bar\sigma}\rangle$. For the uniform ferromagnetic monolayer this is a
rigid shift; writing $\Delta_\uparrow=0$ and $\Delta_\downarrow=\Delta_{\rm ex}$ (the value fixed
by the measured minority gap) reproduces the half-metal. A site-dependent self-consistent version
is only needed for inequivalent edge atoms.

%======================================================================
\section{Landauer currents and Onsager coefficients}

\subsection{Landauer charge and heat currents}
For a two-terminal phase-coherent conductor with transmission $T(E)$ and leads $L,R$ at
electrochemical potentials $\mu_{L,R}$ and temperatures $T_{L,R}$, with Fermi functions
$f_{L,R}(E)=\{\exp[(E-\mu_{L,R})/\kB T_{L,R}]+1\}^{-1}$,
\begin{equation}
I=\frac{e}{h}\int\!\dd E\,T(E)\,[f_L(E)-f_R(E)],\qquad
I_Q=\frac{1}{h}\int\!\dd E\,(E-\mu)\,T(E)\,[f_L(E)-f_R(E)].
\end{equation}
($e>0$; the electron charge is $-e$. $I_Q$ is the heat current carried out of the left lead,
$(E-\mu)$ being the heat per electron.)

\subsection{Linearisation}
Write $\delta\mu=\mu_L-\mu_R$ and $\delta T=T_L-T_R$, small. Since $f$ depends on $E$ only through
$x=(E-\mu)/\kB T$,
\begin{equation}
\frac{\partial f}{\partial\mu}=-\frac{\partial f}{\partial E},\qquad
\frac{\partial f}{\partial T}=-\frac{\partial f}{\partial E}\,\frac{E-\mu}{T},
\end{equation}
which follow from $\partial_\mu x=-1/\kB T$, $\partial_T x=-(E-\mu)/\kB T^2$, $\partial_E x=1/\kB T$.
Hence
\begin{equation}
f_L-f_R\simeq\Big(-\frac{\partial f}{\partial E}\Big)\Big[\delta\mu+\frac{E-\mu}{T}\,\delta T\Big].
\end{equation}
Define the (transport) moments
\begin{equation}
\boxed{\ \Lloc_n(\mu,T)=\int\!\dd E\Big(-\frac{\partial f}{\partial E}\Big)(E-\mu)^n\,T(E)\ }
\qquad(n=0,1,2).
\end{equation}
Then the linear-response currents are
\begin{equation}
I=\frac{e}{h}\Big[\Lloc_0\,\delta\mu+\frac{\Lloc_1}{T}\,\delta T\Big],\qquad
I_Q=\frac{1}{h}\Big[\Lloc_1\,\delta\mu+\frac{\Lloc_2}{T}\,\delta T\Big].
\end{equation}

\subsection{Conductance, Seebeck, thermal conductance}
\paragraph{Electrical conductance} ($\delta T=0$, $\delta\mu=-e\,\Delta V$ where
$\Delta V=V_L-V_R$):
\begin{equation}
I=\frac{e}{h}\Lloc_0(-e\Delta V)=-\frac{e^2}{h}\Lloc_0\,\Delta V
\ \Rightarrow\
\boxed{\,G=\frac{e^2}{h}\,\Lloc_0\,.}
\end{equation}

\paragraph{Seebeck coefficient} (open circuit, $I=0$): setting $I=0$,
$\Lloc_0\delta\mu+\Lloc_1\delta T/T=0\Rightarrow\delta\mu=-(\Lloc_1/\Lloc_0)\delta T/T$. With
$\delta\mu=-e\Delta V$, $\Delta V=(\Lloc_1/e\Lloc_0)\delta T/T$, and $S\equiv-\Delta V/\delta T$:
\begin{equation}
\boxed{\,S=-\frac{1}{eT}\,\frac{\Lloc_1}{\Lloc_0}\,.}
\end{equation}

\paragraph{Electronic thermal conductance} (heat current at $I=0$): substitute
$\delta\mu=-(\Lloc_1/\Lloc_0)\delta T/T$ into $I_Q$,
\begin{equation}
I_Q=\frac{1}{h}\Big[\Lloc_1\Big(-\frac{\Lloc_1}{\Lloc_0}\frac{\delta T}{T}\Big)+\frac{\Lloc_2}{T}\delta T\Big]
=\frac{\delta T}{hT}\Big(\Lloc_2-\frac{\Lloc_1^2}{\Lloc_0}\Big)
\ \Rightarrow\
\boxed{\,\kappa_e=\frac{1}{hT}\Big(\Lloc_2-\frac{\Lloc_1^2}{\Lloc_0}\Big).}
\end{equation}

\subsection{Units}
With energies in eV, $\Lloc_0$ is dimensionless, $\Lloc_1$ has units eV, $\Lloc_2$ eV$^2$.
$G=(e^2/h)\Lloc_0$ is in $e^2/h=3.874\times10^{-5}$~S. For $S$, note $\Lloc_1/\Lloc_0$ in eV
divided by $e$ (the elementary charge) is a voltage: $[\text{eV}]/[e]=\text{V}$. Hence
\begin{equation}
S[\text{V/K}]=-\frac{1}{T[\text K]}\,\frac{\Lloc_1}{\Lloc_0}\Big[\text{eV}\Big]
\quad\Longleftrightarrow\quad
S[\mu\text{V/K}]=-\frac{10^6}{T}\,\frac{\Lloc_1}{\Lloc_0}.
\end{equation}
For $\kappa_e$, convert eV$^2\to$J$^2$ with $Q^2$ ($Q=1.602\times10^{-19}$~C $=$ 1~eV/V):
$\kappa_e[\text{W/K}]=(1/hT)(\Lloc_2-\Lloc_1^2/\Lloc_0)\,Q^2$, with $h$ in J\,s. Equivalently, per
spin $G=G_0\Lloc_0$ with $G_0=Q^2/h$.

\subsection{Electronic figure of merit}
The full $ZT=S^2GT/(\kappa_e+\kappa_{\rm ph})$. Setting $\kappa_{\rm ph}=0$ gives the electronic
bound; substitute the boxed results:
\begin{equation}
S^2GT=\Big(\frac{1}{eT}\frac{\Lloc_1}{\Lloc_0}\Big)^2\frac{e^2}{h}\Lloc_0\,T
=\frac{1}{hT}\frac{\Lloc_1^2}{\Lloc_0},\qquad
\kappa_e=\frac{1}{hT}\Big(\Lloc_2-\frac{\Lloc_1^2}{\Lloc_0}\Big),
\end{equation}
so the $1/hT$ and all $e$ factors cancel:
\begin{equation}
\boxed{\,ZT_e=\frac{S^2GT}{\kappa_e}=\frac{\Lloc_1^2/\Lloc_0}{\Lloc_2-\Lloc_1^2/\Lloc_0}
=\frac{\Lloc_1^2}{\Lloc_0\Lloc_2-\Lloc_1^2}\,.}
\end{equation}
This is manifestly unit-free and diverges when $\Lloc_0\Lloc_2\to\Lloc_1^2$ (a gap, where all
$\Lloc_n\to0$)---the origin of the spurious in-gap divergence that a finite $\kappa_{\rm ph}$
removes.
\begin{quote}\small\itshape
Numerical check: $ZT_e$ computed directly from $\Lloc_1^2/(\Lloc_0\Lloc_2-\Lloc_1^2)$ agreed to
machine precision with $S^2GT/\kappa_e$ evaluated from the separate $G,S,\kappa_e$ routines.
\end{quote}

%======================================================================
\section{Sommerfeld expansion, Mott formula, Wiedemann--Franz}

\subsection{The expansion}
For a smooth $H(E)$, expand about $\mu$ inside $\Lloc[H]=\int(-\partial_E f)H(E)\dd E$. With
$x=(E-\mu)/\kB T$, $(-\partial_E f)\dd E=\tfrac14\,\mathrm{sech}^2(x/2)\,\dd x$, an even, unit-weight
kernel:
\begin{equation}
\int(-\partial_E f)\dd E=1,\quad \int(-\partial_E f)(E-\mu)\dd E=0,\quad
\int(-\partial_E f)(E-\mu)^2\dd E=(\kB T)^2\!\!\int\! x^2\frac{\ee^{x}}{(\ee^{x}+1)^2}\dd x=\frac{\pi^2}{3}(\kB T)^2 .
\end{equation}
(The last integral is a standard result, $\int_{-\infty}^\infty x^2\ee^x/(\ee^x+1)^2\dd x=\pi^2/3$.)
Therefore, to leading order,
\begin{equation}
\Lloc[H]=H(\mu)+\frac{\pi^2}{6}(\kB T)^2\,H''(\mu)+O(T^4).
\end{equation}

\subsection{Moments}
Apply with $H=T(E),\,(E-\mu)T(E),\,(E-\mu)^2T(E)$ (primes $=\dd/\dd E$ at $\mu$):
\begin{equation}
\Lloc_0\simeq T+\frac{\pi^2}{6}(\kB T)^2 T'',\qquad
\Lloc_1\simeq\frac{\pi^2}{3}(\kB T)^2 T',\qquad
\Lloc_2\simeq\frac{\pi^2}{3}(\kB T)^2\,T,
\end{equation}
where $T\equiv T(\mu)$ etc. (For $\Lloc_1$, $H=(E-\mu)T$ has $H(\mu)=0$, $H''(\mu)=2T'$; for
$\Lloc_2$, $H=(E-\mu)^2T$ has $H(\mu)=0$ and the leading term is the second moment $\times\,T$.)

\subsection{Mott formula}
\begin{equation}
S=-\frac{1}{eT}\frac{\Lloc_1}{\Lloc_0}\simeq-\frac{1}{eT}\frac{(\pi^2/3)(\kB T)^2 T'}{T}
=\boxed{\,-\frac{\pi^2}{3}\frac{\kB^2 T}{e}\,\frac{\dd\ln T(E)}{\dd E}\Big|_{\mu}\,.}
\end{equation}
\begin{quote}\small\itshape
Numerical check: for a linear $T(E)=T_0+bE$ at 300~K the code gave $S=-4.397~\mu$V/K, matching
$-\tfrac{\pi^2}{3}\tfrac{\kB}{e}(\kB T)(b/T_0)=-4.40~\mu$V/K.
\end{quote}

\subsection{Wiedemann--Franz law}
In the metallic limit $\Lloc_1^2/\Lloc_0=O\big[(\kB T)^4/T\big]$ is higher order and negligible
against $\Lloc_2=O[(\kB T)^2 T]$, so
\begin{equation}
\kappa_e\simeq\frac{\Lloc_2}{hT}\,Q^2=\frac{1}{hT}\frac{\pi^2}{3}(\kB T)^2 T(\mu)\,Q^2,\qquad
G=\frac{Q^2}{h}T(\mu),
\end{equation}
hence
\begin{equation}
\boxed{\,\frac{\kappa_e}{G\,T}=\frac{\pi^2}{3}\frac{\kB^2}{e^2}=L_0=2.44\times10^{-8}~\text{W}\,\Omega\,\text{K}^{-2}\,,}
\end{equation}
the Lorenz number ($e=Q$).
\begin{quote}\small\itshape
Numerical check: in the metallic region ($\mu-E_F<0$) the code returned $\kappa_e/(GT)$ within
$\sim$2--6\% of $L_0$.
\end{quote}

%======================================================================
\section{Two spin channels (parallel conductors)}

The spin channels are independent and carry currents in parallel across the same $\Delta V,\Delta
T$. From \S3, each channel obeys $I_\sigma=-G_\sigma(\Delta V+S_\sigma\Delta T)$, where we used
$\Lloc_1^\sigma=-eT S_\sigma\Lloc_0^\sigma$ so that
$(e/h)\Lloc_1^\sigma\Delta T/T=-G_\sigma S_\sigma\Delta T$. Charge and heat currents add:
$I=\sum_\sigma I_\sigma$, $I_Q=\sum_\sigma I_{Q\sigma}$.

\paragraph{Conductance.} $\;G=\sum_\sigma G_\sigma$.

\paragraph{Seebeck.} Open circuit $I=\sum_\sigma G_\sigma(\Delta V+S_\sigma\Delta T)=0$ gives
$\Delta V=-\Delta T\,\dfrac{\sum_\sigma G_\sigma S_\sigma}{\sum_\sigma G_\sigma}$, so
\begin{equation}
\boxed{\,S=-\frac{\Delta V}{\Delta T}=\frac{\sum_\sigma G_\sigma S_\sigma}{\sum_\sigma G_\sigma}\,.}
\end{equation}

\paragraph{Thermal conductance (with bipolar term).} The heat current at $I=0$ is
$\kappa=\sum_\sigma\kappa_\sigma$ \emph{plus} a circulating-current (bipolar) correction. Carrying
the algebra, the total electronic thermal conductance at zero \emph{total} charge current is
\begin{equation}
\boxed{\;\kappa_e=\sum_\sigma\kappa_\sigma
+T\Big[\sum_\sigma G_\sigma S_\sigma^2-\frac{\big(\sum_\sigma G_\sigma S_\sigma\big)^2}{\sum_\sigma G_\sigma}\Big].}
\end{equation}
The bracket is non-negative (Cauchy--Schwarz) and represents heat carried by the short-circuit
current that flows between channels when they have different thermopowers.

\paragraph{Magnitude here.} The code (\texttt{combine\_spins}) evaluates the full expression above,
including the bracket. Inside the minority gap $[-3.2,+1.0]$~eV the minority channel carries
nothing ($G_\downarrow=S_\downarrow=\kappa_\downarrow=0$), so the bracket vanishes identically
and the in-window results are unaffected. At the global operating point of the extended model,
$\mu-E_F\simeq+1.05$~eV, \emph{both} spins conduct with very different thermopowers
($S_\downarrow\simeq-150$, $S_\uparrow\simeq+4~\mu$V/K at 300~K), so the bipolar term is active
there and is fully included in every reported number; it penalises exactly the two-channel
counterflow that the minority-edge operating point exploits.

%======================================================================
\section{Ballistic phonon thermal-conductance quantum}

For a single ballistic (fully transmitting) phonon mode of dispersion $\omega(q)$, the heat
current is $I_Q=\frac1{2\pi}\int_0^\infty\dd q\,\frac{\dd\omega}{\dd q}\,\hbar\omega\,[n(\omega,T_L)-n(\omega,T_R)]$
with $n$ the Bose function. Changing variable to $\omega$ and linearising in $\delta T$,
\begin{equation}
\kappa_{\rm ph}^{(1)}=\frac{\hbar}{2\pi}\int_0^\infty\!\dd\omega\,\omega\,\frac{\partial n}{\partial T}
=\frac{\hbar}{2\pi}\int_0^\infty\!\dd\omega\,\omega\,\frac{\hbar\omega}{\kB T^2}\frac{\ee^{\hbar\omega/\kB T}}{(\ee^{\hbar\omega/\kB T}-1)^2}.
\end{equation}
With $u=\hbar\omega/\kB T$, $\displaystyle\int_0^\infty\frac{u^2\ee^u}{(\ee^u-1)^2}\dd u=\frac{\pi^2}{3}$, so
\begin{equation}
\boxed{\,\kappa_0(T)\equiv\kappa_{\rm ph}^{(1)}=\frac{\pi^2\kB^2 T}{3h}=9.464\times10^{-13}\,T~\text{W/K}\,,}
\end{equation}
the universal thermal-conductance quantum (same $\pi^2/3$ as the electronic Lorenz factor, as
required by the Wiedemann--Franz analogy).

\paragraph{Mode counting for a ribbon (production estimate).} $4\kappa_0(T)$ is only the
$T\to0$ limit (four gapless ribbon modes). The production estimate (\texttt{crnte.phonon})
counts modes from the \emph{monolayer phonon dispersion} digitized from the published DFPT
spectrum (Modarresi et al., PRApplied 11, 064015 (2019), Fig.~4(a)): acoustic branches with
$v_{\rm LA}\simeq8.5$, $v_{\rm TA}\simeq4.3$~km/s and zone-edge maxima $4.7$--$10.4$~THz (ZA
quadratic near $\Gamma$), plus three optical branches ($13.7\!\to\!7.2$, $13.2\!\to\!10.8$,
$16.85\!\to\!15.1$~THz from $\Gamma$ to zone edge), each represented by a monotone sine-form
interpolation (whose zone-edge value and slope reproduce the digitized velocities---an internal
consistency check). For an isotropic branch $b$, a ribbon of width $W$ has
$M_b(\omega)=(W/\pi)\,q_b(\omega)$ propagating channels, with $q_b(\omega)$ the iso-frequency
radius; the total $M(\omega)$ is floored at 4 in the acoustic region. Then
$\kappa_{\rm ph}(T)=\frac1{2\pi}\int d\omega\,\hbar\omega\,(\partial n_B/\partial T)M(\omega)$.
This recovers $4\kappa_0 T$ as $T\to0$ (verified numerically to 1\%), grows roughly linearly
with $W$, and saturates above $\sim$500~K when all branches are populated:
$\kappa_{\rm ph}(300~\text{K})=0.6$--$1.7$~nW/K for $W=10$--$32$~\AA. $ZT$ is reported for
$\tfrac12$--$2\times$ this estimate.

%======================================================================
\section{Mean-field edge magnetism and exchange constants}

\subsection{Self-consistent unrestricted Hartree--Fock}
Decoupling the Cr-$d$ Hubbard term $H_U=U\sum_{i\alpha}\hat n_{i\alpha\uparrow}\hat n_{i\alpha\downarrow}$
(\S2.5) gives the spin-dependent on-site potential
$\varepsilon_{i\alpha\sigma}=\varepsilon^0_\alpha+U\langle\hat n_{i\alpha\bar\sigma}\rangle$. In the
Stoner (uniform-over-$d$) form the site exchange splitting is
\begin{equation}
\Delta_i\equiv\varepsilon_{i\downarrow}-\varepsilon_{i\uparrow}=U_{\rm eff}\,m_i,
\qquad m_i=\sum_\alpha\big(n_{i\alpha\uparrow}-n_{i\alpha\downarrow}\big),
\end{equation}
and the occupations, at fixed electron number $N_e$ (charge neutrality fixes $E_F$),
\begin{equation}
n_{i\alpha\sigma}=\frac1{N_k}\sum_{k,n}f\!\big(E_{n\sigma}(k)\big)\,
\big|\langle i\alpha\sigma|\psi_{n\sigma}(k)\rangle\big|^2 .
\end{equation}
Iterating $\{m_i\}\to\Delta_i\to H_\sigma(k)\to\{m_i\}$ to self-consistency yields the moment
profile; edge Cr atoms, with reduced coordination, develop enhanced $m_i$.

\subsection{Bulk calibration}
$U_{\rm eff}$ and $N_e$ are fixed by requiring the uniform (bulk) self-consistent solution to
reproduce the fitted exchange, $U_{\rm eff}=\Delta_{\rm ex}/m_{\rm bulk}$. In the reduced manifold
$m_{\rm bulk}=2.67\,\mu_B$ ($d_{z^2}$ fully polarised, the metallic $d_{xz},d_{yz}$ partially), so
$U_{\rm eff}=3.6/2.67=1.35$~eV and $N_e=4.72$ e/cell; the magnetic solution is a stable fixed point
(converges from any $m_{\rm init}$).

\subsection{Exchange constants: the magnetic force theorem}
The intersite exchange $J_{ij}$ of the mapped Heisenberg model
$H_{\rm spin}=-\sum_{ij}J_{ij}\,\mathbf e_i\!\cdot\!\mathbf e_j$ (unit vectors $\mathbf e$) is the
second-order response of the electronic energy to \emph{infinitesimal} tilts of the moments. A
finite spin flip is unusable: it changes the on-site potential by $\mathcal O(\Delta_i)\sim$~eV,
whose kinetic/hybridisation energy swamps the meV-scale $J$ (this is why the naive total-energy
flip gave a wrong, on-site-dominated answer). The magnetic force theorem removes this: at
self-consistency the first-order change of the total energy under a potential perturbation equals
the change of the band energy at frozen potential, and the double-counting term (a function of
$|m_i|,n_i$ only, unchanged by a rotation) cancels.

\paragraph{LKAG derivation.} Tilt the moment at site $i$ by a small angle $\theta_i$ about an axis
$\perp$ the collinear direction. The local spin potential
$V_i=\tfrac{\Delta_i}{2}\,\hat z\!\cdot\!\boldsymbol\tau$ acquires a transverse piece
$\delta V_i=\tfrac{\Delta_i}{2}\theta_i\,\hat x\!\cdot\!\boldsymbol\tau$ (Pauli matrices
$\boldsymbol\tau$). By the force theorem the second-order change of the grand potential from tilts
at $i,j$ is
\begin{equation}
\delta^2\Omega=-\frac1\pi\,{\rm Im}\!\int^{E_F}\!\!dE\ {\rm Tr}\big[\delta V_i\,G(E)\,\delta V_j\,G(E)\big].
\end{equation}
Since $\hat x\!\cdot\!\boldsymbol\tau$ flips spin ($\uparrow\!\leftrightarrow\!\downarrow$), the spin
trace produces $\tfrac{\theta_i\theta_j}{4}\Delta_i\Delta_j\,{\rm Tr}_L[G^\uparrow_{ij}G^\downarrow_{ji}]$
(plus its conjugate). Matching requires care with the sign: expanding
$H_{\rm spin}=-\sum_{ij}J_{ij}\,\mathbf e_i\!\cdot\!\mathbf e_j$ for tilts about a common axis,
$\mathbf e_i\!\cdot\!\mathbf e_j=\cos(\theta_i-\theta_j)$, so the \emph{cross} term is
$\partial^2 H_{\rm spin}/\partial\theta_i\partial\theta_j=-J_{ij}$; hence
$J_{ij}=-\partial^2\Omega/\partial\theta_i\partial\theta_j$ and
\begin{equation}
\boxed{\ J_{ij}=-\frac1{2\pi}\,{\rm Im}\!\int^{E_F}\!\!dE\ {\rm Tr}_L\!\big[\Delta_i\,G^\uparrow_{ij}(E)\,\Delta_j\,G^\downarrow_{ji}(E)\big]\ }
\end{equation}
with $\Delta_i$ the on-site exchange-splitting matrix ($=U_{\rm eff}m_i$ on the $d$ orbitals) and
${\rm Tr}_L$ over orbitals ($J>0$ ferromagnetic in this convention). An early version of the
code omitted the minus sign; the corrected sign gives $J_1>0$, consistent with the fact that
the ferromagnetic edge is the stable self-consistent solution.

\paragraph{Ribbon evaluation.} The intersite Green's function between an edge Cr and its $n$-th
neighbour along the periodic edge is a 1D $k$-sum at the self-consistent FM Hamiltonian,
\begin{equation}
G^\sigma_{n}(E)=\frac1{N_k}\sum_k\big[(E+i0^+-H_\sigma(k))^{-1}\big]_{ii}\,e^{ikna},
\end{equation}
($ii$ the edge-site $d$-block, $a$ the edge period), so
\begin{equation}
J_n=-\frac{\delta^2}{2\pi}\,{\rm Im}\!\int^{E_F}\!\!dE\ {\rm Tr}_d\big[G^\uparrow_{n}(E)\,G^\downarrow_{-n}(E)\big],
\qquad \delta=U_{\rm eff}\,m_{\rm edge}.
\end{equation}
This yields $|J_1|\gg|J_2|$ ($\sim$ tens and $\sim\!1$~meV for the zigzag edge) with
ferromagnetic $J_1>0$---the
DMRG sign, hierarchy and order of magnitude. Mapping onto a spin-$S$ model uses $J^{\rm Heis}=J/S^2$
with $S=m_{\rm edge}/2$.

\subsection{Spinful domain-wall junction}
For the antiparallel valve with a noncollinear wall the two spin channels mix and the full
spinful Hamiltonian is needed: each orbital block acquires
\begin{equation}
H_{\rm ex}(x)=\tfrac{\Delta}{2}\big[\mathbb 1-\hat n(x)\!\cdot\!\boldsymbol\sigma\big],
\qquad \hat n=(\sin\theta,0,\cos\theta),\qquad
\theta(x)=2\arctan e^{(x-x_0)/\delta},
\end{equation}
(Walker profile, wall width $\lambda=\pi\delta$), which for $\hat n=+\hat z$ reproduces the
collinear convention (majority unshifted, minority $+\Delta$). Hoppings are spin diagonal.
Checks: a uniform rotation ($\theta$ constant everywhere) leaves $T$ exactly invariant
(verified: $T=8.000000$ at $E_F$ for several $\theta$); $\delta\to0$ reproduces the exact
collinear OFF state ($T=0$ at machine precision); $\lambda\gtrsim3$~nm reaches the adiabatic
limit $T\to T_{\rm P}$ exactly (spin tracks the texture). The crossover
($\lambda_{1/2}\simeq9$~\AA\ at $E_F$) is governed by the mistracking scale
$\hbar v_F/\Delta_{\rm ex}\simeq0.4$~\AA\ (Cabrera--Falicov). NOTE: Kwant attaches a lead on
the side its symmetry vector points toward; an early version of the wall builder had the
$\theta{=}0$ lead on the wrong side, creating a spurious collinear-blocked contact -- caught by
the uniform-rotation and adiabatic-limit checks above.

%======================================================================
\section{Summary of results to check}
\begin{itemize}
\item[(i)] $a=\sqrt3\,d$; reciprocal vectors; $f=0$ at $K$ ($\S1$).
\item[(ii)] Planar Slater--Koster selection rule: $p_z\!\leftrightarrow\!(d_{xz},d_{yz})$ only,
$d_{z^2}$ non-bonding ($\S2.3$).
\item[(iii)] $G=\tfrac{e^2}{h}\Lloc_0$, $S=-\tfrac1{eT}\tfrac{\Lloc_1}{\Lloc_0}$,
$\kappa_e=\tfrac1{hT}(\Lloc_2-\tfrac{\Lloc_1^2}{\Lloc_0})$, with units ($\S3$).
\item[(iv)] $ZT_e=\Lloc_1^2/(\Lloc_0\Lloc_2-\Lloc_1^2)$ ($\S3.6$).
\item[(v)] Mott $S=-\tfrac{\pi^2}{3}\tfrac{\kB^2 T}{e}\,\dd\ln T/\dd E$; Wiedemann--Franz
$\kappa_e/GT=\tfrac{\pi^2}{3}(\kB/e)^2$ ($\S4$).
\item[(vi)] Parallel spins: $G=\sum G_\sigma$, $S=\tfrac{\sum G_\sigma S_\sigma}{\sum G_\sigma}$,
and $\kappa_e=\sum\kappa_\sigma+$ bipolar term---the code evaluates the full expression; the
bracket vanishes inside the minority gap and is active (and included) at the minority-edge
operating point ($\S5$).
\item[(vii)] $\kappa_0=\pi^2\kB^2 T/3h$ ($\S6$).
\item[(viii)] Self-consistent mean field $\Delta_i=U_{\rm eff}m_i$ with bulk calibration; LKAG
exchange $J_{ij}=\tfrac1{2\pi}{\rm Im}\!\int^{E_F}\!{\rm Tr}[\Delta_i G^\uparrow_{ij}\Delta_j G^\downarrow_{ji}]$
(a finite spin flip is on-site dominated and must \emph{not} be used) ($\S7$).
\end{itemize}
Items (ii), (iv), (v), the bipolar term and the $J$ hierarchy $|J_1|\gg|J_2|$ were cross-checked
numerically (boxes above / \S7); every formula used in the code is now derived here in full.

\end{document}
